\documentclass[11pt]{article}
\usepackage[a4paper,margin=24mm]{geometry}
\usepackage{amsmath,amssymb,amsthm,booktabs,array}
\newtheorem{theorem}{Theorem}\newtheorem{lemma}[theorem]{Lemma}
\newcommand{\PG}{\operatorname{PG}}\newcommand{\PSL}{\operatorname{PSL}}
\newcommand{\orb}{\mathcal B}\newcommand{\hyp}{\mathcal H}
\title{Ordered butterflies and oval constructions for $\PSL_3(4)$}
\author{L100 finite construction certificate}\date{}
\begin{document}\maketitle
We certify seven concrete constructions in the fixed 22-page Halle preprint
of Gottschalk and Leemans \cite{GL}, using the already certified point action
of $G=\PSL_3(4)$ and its coset geometries. The reference is the preprint,
not an independently collated journal edition. We make the butterfly
convention and the otherwise unspecified oval classes explicit. This is
a construction component, not a new classification or an assertion that
all remaining structural names in the source have been proved.

\paragraph{Acquired finite premises.}
The coordinate model is the plane over
$\mathbb F_2[t]/(t^2+t+1)$, with an explicit bijection from its normalized
coordinate vectors to the retained 21-point permutation model. Its saved
20160 permutations and two generators are the same literal group as the
coset catalogue. The acquired point/line/hyperoval certificate gives the
21 lines and three 56-element hyperoval orbits. The acquired Baer certificate
gives the three 120-element orbits $\orb_0,\orb_1,\orb_2$ of Fano subplanes.
Completeness of these 360 subplanes follows from the frame count
$21\cdot20\cdot16\cdot9/(7\cdot6\cdot4)=360$, with a unique embedded
Fano plane through a frame. These literal inputs, their proofs and their
distinct reviews are acquired without repeating a group or subgroup census.

\paragraph{An explicit convention.}
An ordered butterfly in this note is an ordered pair of unordered poles
$b=(\{B_1,B_2\},\{C_1,C_2\})$, consisting of four distinct subplanes,
such that, for a point $p$ and a line $z$,
\[
 B_i\cap C_j=\{p\}\quad(1\leq i,j\leq2),\qquad
 T_+=B_1\cap B_2,\quad T_-=C_1\cap C_2,
 \quad |T_+|=|T_-|=3,\quad T_+\cup T_-=z.                 \tag{1}
\]
This all-cross-intersections convention is stronger than the two displayed
cross-intersection conditions on source page4; the distinction is addressed
on the last page. Write $U_+=B_1\cup B_2$, $U_-=C_1\cup C_2$ and
\[
 Q_+=T_+\setminus\{p\},\qquad
 L_\pm=\{l:p\in l\subseteq U_\pm\},\qquad
 H_+=\{(B_i\setminus z)\cup(T_-\setminus\{p\}):i=1,2\},
\]
with $H_-$ defined by interchanging the poles. Thus a line ``contained in
a pole'' means containment in the union of its two point sets. The central
line is not such a line: its intersection with that union has only three points.

\begin{lemma}
For each $\orb_i$, (1) gives exactly1260 ordered butterflies, one orbit of
$G$, with stabilizer $D_8\times C_2$. Both $L_+$ and $L_-$ have two lines.
The four members of $H_+\cup H_-$ are distinct hyperovals in one orbit
$\hyp_{e(i)}$, and $\orb_i$ meets every hyperoval in this orbit evenly.
The ordered positive pole determines the butterfly uniquely. Its unordered
four-subplane set is also determined by an opposite pair of subplanes or by
one of its subplanes together with $(p,z)$.
\end{lemma}
\begin{proof}
The finite certificate has a complete, small domain. Test all
$\binom{360}{2}=64620$ subplane pairs. Retain a pair when its intersection
is a collinear triple $T$. There are3780 such pairs. For each, its line $z$
is unique. For each of the three choices $p\in T$, test all retained pairs
whose intersection is $(z\setminus T)\cup\{p\}$, and impose all four
equalities in (1). This exhausts every possible opposite pole, and gives
exactly one. The resulting objects have1260 members in each Baer orbit.
Closure and transitivity are checked using the two actual group generators.
The alternative uniqueness assertions are checked by grouping the full list
by the indicated exact subplane/flag keys, never by numerical invariants.

For a representative in each orbit, test all saved group permutations for
stabilization of the ordered poles. The 16 retained permutations are verified
closed. The certificate exhibits $a,b,z$ with $|a|=4$, $|b|=|z|=2$,
$bab=a^{-1}$ and $[a,z]=[b,z]=1$; its sixteen words
$a^ib^jz^k$, $0\leq i<4$, $0\leq j,k<2$, are distinct and exhaust the
stabilizer. This proves its asserted structure, not merely its order.
The source's auxiliary subplane-stabilizer transitivity is checked on all
21 partners of one subplane in each orbit.

Each asserted hyperoval is checked against the complete acquired list.
All $168\binom64=2520$ four-subsets of these hyperovals are distinct;
therefore the associated quadrangle has a unique such completion.
Containment tests give the two pole-line sets, the even-intersection
assertion and all stated uniqueness keys. These are finite set tests on the
complete domains just described; the replay recomputes them.
\end{proof}

\paragraph{The six butterfly geometries.}
For $B\in\orb_i$, set
$R(B,b)\iff T_+\subseteq B\text{ and }B\notin\{B_1,B_2\}$.
The following table states every incidence relation between distinct types.
Here $h,h'$ range over two tagged copies of the same class
$\hyp_{e(i)}$, $x$ is a point, $l$ a line, and $b$ an ordered butterfly
from $\orb_i$. The ordered types and the six entries are respectively the
relations12,13,14,23,24,34. Equal-type incidence is equality.
\begin{center}\small
\renewcommand{\arraystretch}{1.25}
\begin{tabular}{ccl}
\toprule Source&Types&Relations in the stated order\\\midrule
4.2&$l,x,h,b$&$x\notin l;\ l\cap h=\varnothing;\ l\in L_+;\ x\in h;\ x\in Q_+;\ h\in H_-$\\
4.3&$h,h',B,b$&$h\cap h'=\varnothing;\ |h\cap B|=4;\ h\in H_-;$\\
&&$h'\cap B=\varnothing;\ h'\in H_+;\ R(B,b)$\\
4.4&$h,h',x,b$&$h\cap h'=\varnothing;\ x\in h;\ h\in H_-;\ x\notin h';\ h'\in H_+;\ x\in Q_+$\\
4.5&$l,x,h,b$&$x\notin l;\ l\cap h\ne\varnothing;\ l\in L_-;\ x\in h;\ x\in Q_+;\ h\in H_-$\\
4.6&$x,h,B,b$&$x\in h;\ x\in B;\ x\in Q_+;\ |h\cap B|=4;\ h\in H_-;\ R(B,b)$\\
4.7&$l,x,B,b$&$x\notin l;\ |l\cap B|=3;\ l\in L_+;\ x\in B;\ x\in Q_+;\ B\in\{B_1,B_2\}$\\
\bottomrule
\end{tabular}\end{center}
These are the native signs on pages16--18. In particular4.5 uses type-one
hyperovals from4.4 and the negative pole for its line relation.

\paragraph{Oval and dual-oval classes.}
Use the coordinate polarity $[a]\mapsto\{[v]:a\cdot v=0\}$ with its
literal transport to the saved point and line labels. It induces a permutation
$d$ of the three Baer classes. For each hyperoval orbit, delete each of the
six points from each hyperoval. This gives336 distinct ovals in one orbit;
the deleted point is its nucleus, since the five tangent lines are checked
to pass through it. Polarity gives the corresponding dual ovals, with a
nucleus line. For each ordered pair $i\ne j$, choose
\[
 B\in\orb_i,\quad C\in\orb_j,\quad
 O\in\operatorname{Ovals}(\hyp_{e(i)}),\quad
 D\in\operatorname{DualOvals}(\hyp_{e(d(j))}).             \tag{2}
\]
In the saved class ordering, $e=(0,1,2)$ and $d=(0,2,1)$.
The indices are part of the explicit model, not intrinsic names.
Writing $n(O)$ for a point and $n(D)$ for a line, the source4.11 relations are
\begin{align*}
 B\sim C&\iff B\cap C=\varnothing,&
 B\sim O&\iff |B\cap O|=4,\\
 B\sim D&\iff |B\cap n(D)|=3\ \text{and }|B\cap l|=1\ (l\in D),&
 C\sim O&\iff n(O)\in C\text{ and }C\cap O=\varnothing,\\
 C\sim D&\iff \#\{l\in D:|C\cap l|=3\}=4,\\
 O\sim D&\iff |O\cap n(D)|=2\ \text{and }(O\setminus n(D))\cap l=\varnothing\ (l\in D).
\end{align*}
All six choices in (2) are certified. The source names an oval class and
a dual-oval class without specifying their relative identification; (2)
supplies a successful precise choice, not a claim about arbitrary choices.

\begin{theorem}
For each of the three Baer classes, the six butterfly recipes above are
equivariantly isomorphic, up to an explicitly given automorphism of $G$
and type permutation, to source geometries4.2--4.7. All six choices in (2)
similarly give geometry4.11. Consequently these constructions are firm,
residually connected, satisfy $(IP)_2$, and have a flag-transitive RWPRI
$G$-action, with the local action qualifications stated below.
\end{theorem}
\begin{proof}
We use the accepted coset catalogue, not a new recognition by orders or
diagrams. The source-to-catalogue indices are respectively
$24,19,22,23,18,17,25$. For each of the24 maps, the certificate records
a type permutation, one base object fixed by each actual parabolic, an
explicit group automorphism $\alpha$, and every object image. The map on
right cosets is
$P_tg\mapsto X_t^{\alpha(g)}$.
Every one of the20160 group entries in each of the four coset partitions
is tested: all members of a coset give the recorded object. The recorded
maps are bijections on each type. For every cross-type pair, incidence
and nonincidence agree exactly; the two generator actions agree on every
object. Thus these are actual typed incidence and action isomorphisms.

The automorphisms used are explicit point normalizers and coordinate duality.
Point normalizers are checked on the whole saved group; duality is checked
on every generator edge of its Cayley graph. The actual coordinate-to-retained
label bijection is always used. None of these maps is inferred from a name,
group order, or an unlabelled graph comparison.

The cited geometric properties, including extension of every partial flag,
are acquired for these exact coset geometries. Isomorphisms transport all
flags, chambers, residues and stabilizer actions, proving the assertions
for the concrete constructions. Local double transitivity is not used as
a recognition filter: in the displayed source type order,4.2 has flags
$(1,1,1,0)$,4.11 has $(1,1,0,0)$, and4.3--4.7 have four true flags.
These follow from the acquired literal panel actions and the recorded type maps.
\end{proof}

\paragraph{A narrow qualification to the printed butterfly description.}
On page4 the source displays only
$B_1\cap B'_1=B_2\cap B'_2=\{p\}$, together with the triple/line
condition. Those two displayed requirements alone do not imply the next
sentence's uniqueness from $(B_1,B_2)$. Here is a direct witness in the
retained point labels. Fix
\[
 B_1=\{1,2,3,6,7,10,11\},\qquad B_2=\{1,4,5,8,9,10,11\}.
\]
Two different completions, in the same Baer orbit as $B_1,B_2$, are
\[
\begin{array}{c|l|l}
p&B'_1&B'_2\\\hline
1&\{1,12,13,15,16,18,20\}&\{1,12,13,14,17,19,21\}\\
10&\{4,8,10,12,13,16,19\}&\{3,7,10,12,13,14,18\}.
\end{array}
\]
Both have central line $\{1,10,11,12,13\}$ and satisfy exactly the two
printed cross equalities. In the first, all four cross intersections have
size1; in the second their sizes, by rows, are $(1,3)$ and $(3,1)$.
All six point sets are checked against the complete saved subplane list;
actual permutations of $G$ mapping $B_1$ to each one are supplied. Hence
neither same-orbit restrictions nor forgetting the pole ordering remove
the ambiguity. This establishes insufficiency of the displayed conditions,
not a claim that the conventional butterflies in the source's cited
literature have that broader meaning. We do not inspect or contradict that
external convention here. Convention(1) makes the intended constructions
and the stabilizer assertion completely explicit. No principal
classification is refuted by this ancillary qualification.

\paragraph{Reproducibility and attribution.}
The packet retains the initial failed uniqueness assertion under the weaker
displayed rule, its diagnosis, and the successful successor. Default replay
checks finite object completeness and the24 supplied maps; it does not
search for maps, rerun a subgroup census or recompute graph automorphism
groups. The names and recipes are due to Gottschalk--Leemans and their cited
antecedents. This certificate gives literal construction joins and an explicit
scope qualification, with no novelty, priority, measurement or whole-source claim.
\begin{thebibliography}{9}
\bibitem{GL} H. Gottschalk and D. Leemans, \emph{Geometries for the group
$\PSL(3,4)$}, Halle preprint99--30, fixed22-page source, pp.4--5,16--19.
\end{thebibliography}
\end{document}
